\section{Methodology}
This study applies a Multi-Agent Reinforcement Learning approach to network routing. Each node is represented as an agent that observes local link state and forwards the traffic transiting it.

\subsection{Problem Formulation}
We model a backbone as a directed graph $G=(\mathcal{N},\mathcal{E})$, where
$\mathcal{N}$ is the set of $N$ nodes (routers) and
$\mathcal{E}\subseteq\mathcal{N}\times\mathcal{N}$ the set of directed links. Each
link $e=(i,j)$ has a capacity $c_e>0$ and a propagation delay $\delta_e$, and the
adjacency matrix $A$ gives $A_{ij}=1$ if $(i,j)\in\mathcal{E}$ and $0$ otherwise.
The topology is fixed within an episode but varies between them, since one policy
is trained over a distribution of networks rather than fitted to a single graph
(Section~\ref{sec:method-protocol}).

Traffic is represented as a demand matrix decomposed into flows$\mathcal{F}=\{f=(s_f,d_f,r_f)\}$, where flow $f$requests requests rate $r_f$ from $s_f$ to
$d_f$. During an episode, one matrix is routed and stays the same for the whole period. Training uses synthetic matrices, while evaluation uses real measured demand (Section~\ref{sec:setup-instances}).
A routing assigns each flow a single path $p_f$, giving link loads and utilizations
\begin{equation}\label{eq:load}
\ell_e=\sum_{f\in\mathcal{F}} r_f\,\mathbb{1}\!\left[e\in p_f\right],
\qquad
u_e=100\,\frac{\ell_e}{c_e}\quad[\%],
\end{equation}
and the bottleneck (maximum link) utilization
\begin{equation}\label{eq:bottleneck}
U=\max_{e\in\mathcal{E}} u_e .
\end{equation}
When $U>100\%$ a link is overloaded and the excess appears at the packet level as
loss and queueing delay. The traffic-engineering objective is therefore to minimize
the bottleneck $U$ while preferring short paths, penalising each \emph{detour hop}
of a chosen path, those that do not reduce the OSPF-cost distance to the
destination, at a weight $\beta\ge0$. Counting detours rather than excess hops
matters where capacities are heterogeneous, since the cheapest path in hops and in
OSPF cost then differ.

We formulate routing as a multi-agent Markov Decision Process
$(\mathcal{N},\mathcal{S},\{\mathcal{A}_i\},\allowbreak P,R,\{\mathcal{O}_i\},\gamma)$
under centralized training with decentralized execution (CTDE), with transitions $P$
given by \eqref{eq:load}--\eqref{eq:bottleneck} and $\gamma$ the discount factor:
\begin{itemize}
\item $\mathcal{N}$: set of agents, one per node, each deciding how flows
transiting it are forwarded.
\item $\mathcal{S}$: global state, comprising the link utilizations $\{u_e\}$, the
adjacency $A$, and the metadata $(s_f,d_f,r_f)$ of the flow being routed.
\item $\mathcal{A}_i(t)$: which neighbor to send the flow to next. Only the hops
that \eqref{eq:mask} admits may be chosen, so the agent can never revisit a node or
stray too far from the destination. The centralized reference instead picks one of
$k$ whole candidate paths, $\mathrm{Discrete}(k)$.
\item $R$: one reward shared by every agent, paying for any reduction in the
busiest link's load and charging $\beta$ for each hop that does not move closer to
the destination, as defined in \eqref{eq:reward}.
\item $\mathcal{O}_i(t)$: local observation, a subset of $\mathcal{S}$ holding the
utilizations of the links incident to $i$, its local topology, and $(s_f,d_f)$.
\end{itemize}
Each agent learns a policy $\pi_i:\mathcal{O}_i\to\mathcal{A}_i$ that aims to maximize the shared return. A centralized critic reads $\mathcal{S}$ during training because each agent only sees its own links, and all agents learn at the same time. Then, the critic is removed after training, leaving the execution to rely on $\mathcal{O}_i$ alone.

Loop freedom is enforced structurally rather than learned. The set of allowed next hops is defined as \begin{equation}\label{eq:mask}
\mathcal{A}^{\mathrm{ok}}_i(d)=\bigl\{\,j\in\mathcal{A}_i \;:\;
j\notin\mathcal{V}_f,\;\; h_d(j)<h_d(i)+\sigma,\;\;
\omega_t+w_{ij}+h_d(j)\le h^{\star}_f+\sigma_{\max}\bigr\},
\end{equation} 
where $h_d(i)$ is the distance from $i$ to $d$ under the OSPF cost metric $w_e=\mathrm{refBW}/c_e$. $\mathcal{V}_f$ is the set of nodes already visited,  $\omega_t$ is the cost accumulated so far, $\sigma$ is the limitation of how far a single hop can move from the destination, and $\sigma_{\max}$  is the extra cost allowed for the completed path. Both limits are measured in cost units rather than hops, so detours are judged by capacity. This set is used as a binary mask in the policy. If the set is empty, the agent chooses the unvisited neighbor with the smallest $h_d$.

With $U_t$ the bottleneck \eqref{eq:bottleneck} after the $t$-th hop and $U_0=0$,
the reward is
\begin{equation}\label{eq:reward}
r_t=-\,\frac{100\,(U_t-U_{t-1})}{U^{\mathrm{OSPF}}}\;-\;\beta\,\mathbb{1}\!\left[\text{hop
$t$ is a detour}\right].
\end{equation}
Normalising by $U^{\mathrm{OSPF}}$, the bottleneck OSPF reaches on the same matrix,
keeps networks at different utilizations comparable to a single shared policy, and
the factor $100$ puts the term in percentage points: at $\beta=0.5$ one detour hop
costs half a point of OSPF's bottleneck.

Because each hop is charged only the \emph{change} in the bottleneck and the episode
starts empty, the intermediate values cancel:
\begin{equation}\label{eq:telescope}
\sum_{t} r_t=-\,\frac{100\,U_T}{U^{\mathrm{OSPF}}}\;-\;\beta\sum_{t}
\mathbb{1}\!\left[\text{hop $t$ is a detour}\right].
\end{equation}
Maximizing the return therefore minimizes the objective above exactly, no terminal
bonus is required, and the return is directly readable: $-100$ is OSPF's bottleneck
reached with no detours. Both learners optimise this reward.

Five assumptions bound the formulation and, with it, the scope of the results that
follow. Routing is \emph{unsplittable}: each flow follows the single path of
\eqref{eq:load} rather than being divided across several, which is what allows a
path to be the unit of decision but forgoes the finer balancing that fractional
splitting permits. The topology is \emph{static} within an episode, with capacities
and propagation delays held constant and no link or node failures, so the results
speak to congestion rather than to resilience. Decisions are taken \emph{per demand
matrix} rather than per packet and are stable between matrices; the resulting paths
are source-specific, in that two flows toward the same destination may be routed
differently, which plain destination-based IP forwarding cannot express and a
deployment would have to realise through MPLS or segment routing. Traffic is
\emph{open-loop}, offered as constant-rate UDP, so sources do not back off and the
loss reported is queue overflow alone rather than the milder loss a
congestion-controlled sender would produce. Finally, the analytical model of \eqref{eq:load}
assumes offered demand is carried, an assumption that fails precisely once a link
saturates; it is therefore used only to train, and every reported number is measured
in \mbox{ns-3}.

\subsection{System Architecture}
The policy is a graph neural network (GNN) shared by every node agent. This GNN represents the network state through its nodes and links \citep{Rusek_2020}. Then, we use a custom multi-agent Proximal Policy Optimization (MAPPO) algorithm under centralized training with decentralized execution\citep{marl-book} to train the policy.

Every node $i$ carries a state vector
\begin{equation}\label{eq:nodefeat}
o_i(t)=\bigl[\,\mathbb{1}[i{=}c_t],\;\mathbb{1}[i{=}d_f],\;
\tfrac{h_{d_f}(i)}{h_{\max}},\;
\tfrac{1}{100}\!\max_{e\in\mathrm{out}(i)}\! u_e,\;
\tfrac{1}{100}\!\min_{e\in\mathrm{out}(i)}\!(100{-}u_e),\;
\tfrac{\deg(i)}{\Delta}\,\bigr]\in\mathbb{R}^{6},
\end{equation}
where $c_t$ is the node holding the flow and $\Delta$ the padding bound on degree.
The vector contains \emph{no node identity and no position}: it describes a node
only by its role in the current decision and by the congestion around it, which is
what allows one parameter set to serve a network of any size.

Node embeddings come from an encoder followed by $L=3$ rounds of neighborhood
aggregation with a residual connection,
\begin{equation}\label{eq:gnn}
h^{(0)}_i=\mathrm{ReLU}\bigl(W_{\mathrm{enc}}\,o_i(t)\bigr),
\qquad
h^{(l)}_i=h^{(l-1)}_i+\mathrm{ReLU}\Bigl(W^{(l)}\bigl[\,h^{(l-1)}_i
\;\Vert\;
\bar{A}_i h^{(l-1)}\,\bigr]\Bigr),
\end{equation}
where $\bar{A}=\Lambda^{-1}(A+I)$ is the row-normalized adjacency with self-loops,
so $\bar{A}_i h^{(l-1)}$ is the mean embedding of $i$ and its neighbors, and
$W_{\mathrm{enc}},W^{(l)}$ are learnable. After $L$ rounds, $h^{(L)}_i$ summarizes
the state of a neighborhood of radius $L$ around node $i$, giving each agent context
beyond the links it observes directly without requiring full observability.

The agent at node $c_t$ scores each neighbor individually rather than emitting a
vector over a fixed action set:
\begin{equation}\label{eq:policy}
\eta_j=\mathrm{MLP}_\theta\bigl[\,h^{(L)}_{c_t}\Vert h^{(L)}_{j}\Vert
\ell_{c_tj}\,\bigr],
\qquad
\pi_\theta(j\mid o,m)=\frac{m_j\exp\eta_j}{\sum_{j'}m_{j'}\exp\eta_{j'}},
\end{equation}
where $\ell_{c_tj}\in\mathbb{R}^{4}$ represents the utilization of the link to $j$, the bottleneck for that link, the remaining headroom, and the normalized distance from node $j$ to the destination. The variable m indicates the mask defined in \eqref{eq:mask}. As a result, unacceptable hops are assigned zero probability. This approach ensures every sampled action avoids loops by design. Since $\mathrm{MLP}_\theta$ is used for each neighbor and shared across all neighbors, the policy works for any node degree. On the other hand, an output layer matched to the action set would not generalize across different topologies.

A centralized critic $V_\phi(s_t)$ reads the global state, the utilization of
every link, the current bottleneck, and the normalized rate of the flow being
routed. This critic serves only to compute advantages during training and is discarded at execution. Consequently, deployment requires the node features, \eqref{eq:gnn}, and the scoring network. 

\subsection{Training and Evaluation Protocol}
\label{sec:method-protocol}
The policy is trained using the analytical model described in \eqref{eq:load} and \eqref{eq:bottleneck}, which costs milliseconds per episode. Next, we evaluate the model with packet-level simulation in ns-3. This simulation assumes that all offered demand is carried, but this assumption breaks down when a link becomes saturated. As a result, the simulator is not part of the gradient path, but all reported results are based on its output.

An episode is a complete routing of one demand matrix, its flows ordered by decreasing rate and
routed hop by hop. Training maximizes the expected discounted return \eqref{eq:telescope} over
the training distribution $\mathcal{D}_{\mathrm{train}}$. That expectation runs over a
distribution of \emph{networks} $G$ as well as of demand matrices,
which is what makes the trained object a topology-agnostic policy rather than a routing table:
$\mathcal{D}_{\mathrm{train}}$ ranges over topologies none of which is an evaluation network, so
the separation is \emph{topological} and the generalization measured is zero-shot transfer to
unseen graphs. Instances and traffic are given in
Sections~\ref{sec:setup-instances}.

Each iteration collects $T$ decision steps under $\pi_{\theta_{\mathrm{old}}}$ old and resets to a new matrix at the end of every episode. Advantages $\tilde{A}_t$ are calculated with generalized advantage estimation with a bootstrap at the truncation point. $\tilde{A}_t$  are also standardized across the batch, and the batch is reused for $E$ epochs of minibatch descent on 
\begin{equation}\label{eq:loss}
\mathcal{L}=-\mathbb{E}_t\bigl[\min(\rho_t\tilde{A}_t,\;
\mathrm{clip}(\rho_t,1{\pm}\varepsilon)\tilde{A}_t)\bigr]
+c_v\,\mathbb{E}_t\bigl[(V_\phi(s_t)-\hat{V}_t)^2\bigr]
-c_e\,\mathbb{E}_t\bigl[\mathbb{H}[\pi_\theta]\bigr],
\end{equation}
with $\rho_t=\pi_\theta(a_t\mid o_t,m_t)/\pi_{\theta_{\mathrm{old}}}(a_t\mid o_t,m_t)$. Because the mask enters $\pi_\theta$ through \eqref{eq:policy}, the entropy term only explores legal next hops. Global gradient-norm clipping is used for updates. This is necessary because the reward is based on the maximum over links, and a single hop that creates a new bottleneck can cause a large sparse gradient. All router agents update the same $\theta$, so each iteration uses $T$ decisions instead of $T/N$ per-agent trajectories. We use one configuration in  Table ~\ref{tab:hyper} for every training topology and apply it to every evaluation network. The policy is topology-agnostic, so tuning it for each network would go against its purpose.

Training progress is monitored by analyzing held-out matrices under the deterministic policy $a = \arg\max_a \pi_\theta(a\mid o,m)$, measuring the average bottleneck improvement over OSPF. Subsequently, we use ns-3 in export-simulate form for the final evaluation. The deterministic rollout produces a specific node path for each demand, while OSPF provides its shortest path.  Both are installed as static per-hop routes, so the only difference between runs is the routing decision. Details on the simulator setup and metric definitions are given in Sections ~\ref{sec:setup-ns3} and~\ref{sec:setup-metrics}.
